\section{Localization and Cultural Identity}

The conceptual core of \plat{} is the decision to localize not only
documentation or diagnostics, but the programming surface itself. The language
does this through keyword choice, built-in names, runtime type names,
localized diagnostics, examples, and documentation structure.

\subsection{Keywords and Names}

The core vocabulary includes declaration words such as \code{loat} and
\code{funksie}; control-flow words such as \code{es}, \code{angesj},
\code{zolang}, \code{veur}, \code{enj}, \code{trok}, \code{aafbraeke}, and
\code{euversjlaon}; and literal words such as \code{niks}, \code{woar}, and
\code{neetwoar}. The minimal standard library currently exposes
\code{aafdrokke} for printing and \code{invuier} for input. The protected-name
set also includes runtime type names such as \code{nommer}, \code{teks}, and
\code{portefeuil}. These choices make Limburgian visible at the level where many
languages default to English.

The naming layer is not merely decorative. It changes what is repeated in
tutorials, error messages, tests, and examples. A learner writing
\code{funksie fib(n):} and \code{trok n} encounters programming concepts
through a regional-language vocabulary. This does not remove the need to learn
general programming ideas, but it changes the cultural frame through which
those ideas are introduced.

\subsection{Diagnostics and Documentation}

The implementation includes localized diagnostic templates for English and
Limburgian. The command-line option \code{--lang li} selects Limburgian
diagnostics, while \code{--lang en} selects English. The implementation also
accepts more specific Limburgian tags such as \code{li-sit} and
\code{li-x-sittard}. Diagnostics use stable identifiers and localized message
templates; for example, unknown variables, invalid table keys, wrong arity,
breaks outside loops, and returns outside functions are rendered through the
selected diagnostic language.

The documentation tree mirrors this concern by providing both English and
Limburgian paths. The example tree is similarly split between English-commented
and Limburgian-commented examples. This matters because localization is rarely
only a lexical substitution problem. A language also needs examples, tutorial
phrasing, diagnostic tone, and documentation conventions that make the chosen
identity feel coherent.

\subsection{Dialect Selection}

The design notes state that \plat{} prefers the Sittard dialect as its
foundation where possible. That choice gives the project a concrete linguistic
anchor, but it also raises questions. Limburgian is not a single homogeneous
surface. Different speakers may prefer different spellings, dialect forms, or
degrees of standardization. A programming language that adopts one variety
therefore needs to explain whether it is designing for a local community, a
broader regional identity, a pedagogical compromise, or an artifact authored
from one particular linguistic position.

% TODO: Add evidence from Limburgian speakers, teachers, learners, or language
% organizations on whether the selected terms feel natural, legible, and
% respectful.

\subsection{Authenticity and ASCII Compatibility}

The repository explicitly discusses the tension between authentic written
Limburgian and practical typing. Standard Limburgish spelling often uses
accents and diacritics, but those characters can be awkward on common keyboard
layouts and toolchains. \plat{} therefore keeps official keywords and standard
built-in names within plain ASCII while allowing extended characters in source
files where naturally useful, especially strings and comments.

This is a pragmatic compromise. ASCII-only keywords lower friction for typing,
version control, shells, tutorials, and cross-platform examples. At the same
time, ASCII can flatten orthographic distinctions that matter to speakers.
The design question is not whether authenticity or usability should always
win, but how a language should document and revisit the compromise.

\subsection{Beyond Novelty}

Localized syntax can be dismissed as a gimmick if it is treated as a mere
translation of English keywords. \plat{} avoids that reduction when its
linguistic choices are tied to a broader design program: localized diagnostics,
Limburgian examples, dialect reflection, and explicit discussion of typability
and preservation. The language remains playful, but its playfulness is attached
to a serious question: how can programming systems give regional languages a
visible digital role without pretending that keyword translation alone solves
accessibility, education, or preservation?

% TODO: Add a user study or qualitative feedback section before making strong
% claims about accessibility, identity, or cultural impact.
